\section{Week 1}
\sol{1.1}{At $x=4$, $2(x+3)=2(4+3)=14$, whereas $2x+3=2\cdot4+3=11$. In the first expression we add three before doubling, so the added three is also doubled. In the second we double only the original input and then add three. The difference is three, not a contradiction: these are different functions of the same input.}
\sol{1.2}{The conclusion asks for an odd sum. By the definition, we should aim to write that sum as twice an integer plus one.

Name the even input $a$ and the odd input $b$. Their definitions supply integers $k,\ell$ with $a=2k$ and $b=2\ell+1$. The witnesses may be different because the two inputs were not required to be equal. Substitute and regroup:
\begin{align*}
 a+b&=2k+(2\ell+1)\\
 &=2k+2\ell+1\\
 &=2(k+\ell)+1.
\end{align*}
The number $k+\ell$ is an integer, so the final expression meets the definition of odd. It is the requested witness. For example, $a=4,b=-3$ have witnesses $k=2,\ell=-2$; the sum is one, witnessed by $k+\ell=0$. The example illustrates the proof but does not replace its arbitrary-input argument.}
\sol{1.3}{The proposed calculation uses the same odd number twice. Squaring $2k+1$ therefore covers pairs of equal inputs only. It does not cover, for example, the product of three and five. Correct algebra can still have insufficient scope.

For the general claim, name the inputs independently: $a=2k+1$ and $b=2\ell+1$, with integer $k,\ell$. Then distribute one bracket at a time:
\begin{align*}
 ab&=(2k+1)(2\ell+1)\\
 &=2k(2\ell+1)+(2\ell+1)\\
 &=4k\ell+2k+2\ell+1\\
 &=2(2k\ell+k+\ell)+1.
\end{align*}
The expression in parentheses is an integer because products and sums of integers are integers. This gives the defining form of an odd integer and proves the claim for every allowed pair. Setting $k=\ell$ now recovers the square case as one consequence, not as a substitute for the general proof.}
\sol{1.4}{The sum is
\[
 (2a_1+1)+(2a_2+1)+(2a_3+1)
 =(2\cdot2+1)+(2(-1)+1)+(2\cdot4+1)
 =5-1+9=13.
\]
There are three terms because the index runs from one through three, not because the upper limit is a multiplicative factor.}
\sol{1.5}{A zero product requires $x=0$ or $x-2=0$. Thus the only candidates are $0$ and $2$. Substitution verifies both. Dividing by $x$ discards the valid candidate $0$, because division is permitted only when $x\ne0$. A division-based proof is complete only after separately checking the excluded case.}
\sol{1.6}{Use
\[
 f(n)=(n-1)(n-2)\cdots(n-10).
\]
At an integer $n$ from one through ten, one factor is zero. At $n=11$, all factors are positive, giving $10\cdot9\cdots1=3\,628\,800$. The finite observations do not determine the behavior of this formula outside the tested inputs.}
